% rhythm: article display math: a long last line into an unnumbered display and out of it, then the same around a numbered equation
% boundary: par-display Birch span
% boundary: display-par Cedar span
% boundary: par-display Dogwood span
% boundary: display-par Elm span
\documentclass{article}
\usepackage{amsmath}
\usepackage{fontspec}
\setmainfont{OpenSans-Regular.ttf}[Path=../corpus/fonts/, BoldFont=OpenSans-Bold.ttf, ItalicFont=OpenSans-Italic.ttf, BoldItalicFont=OpenSans-BoldItalic.ttf]
\usepackage{unicode-math}
\setmathfont{FiraMath-Regular.otf}[Path=../corpus/fonts/]
\begin{document}
Alder sets one line that reaches past the display.
\[ a + b = \text{Birch} \]
Cedar follows the unnumbered display with one line too.
\begin{equation} c + d = \text{Dogwood} \end{equation}
Elm closes the page after the numbered equation.
\end{document}
